\textbf{Cálculos energéticos del sistema}

En esta parte se realizan los cálculos para determinar cuánta energía eléctrica es capaz de consumir el dispositivo y, de esta manera, realizar una aproximación del costo económico que representaría.

\textbf{Calculo máximo de energía.}

En este cálculo se tiene en cuenta que todos los subsistemas están encendidos al mismo tiempo; si bien es una opción poco probable, con este valor se puede determinar qué tipo de instalación se requiere para ponerlo en marcha.

Para cada uno de los elementos eléctricos del sistema mecatrónico,se tienen los siguientes valores de potencia eléctrica (Tabla \ref{tab:Electrical_values}).


\begin{table}
    \centering
\caption{Valores de potencia de los elementos eléctricos del sistema}
\label{tab:Electrical_values}
    \begin{tabular}{|m{2.5cm}|m{2.5cm}|m{2.5cm} |m{2.5cm}|} \hline 
         Elemento&  Potencia Activa nominal (W)&Factor de Potencia FS &Potencia aparente (VA)\\ \hline 
         Motor de trituración&  1397&1.35 &1034.815\\ \hline 
         Motor de aireación&  180&1.35 &133.333\\ \hline 
         Circuito electrónico&  20&1 &20\\ \hline 
         Resistencia eléctrica&  1000&1 &1000 \\ \hline 
         Ventilador&  12&1 &12\\ \hline 
         Fuente de 24 V&  120&1 &120 \\\hline
    \end{tabular}
    
    
\end{table}

Para encontrar la corriente máxima $I_{MAX}$ que será necesaria para alimentar el circuito, es necesario encontrar la potencia aparente total del sistema $S_T$  y esta es dividida por el valor de la tensión de alimentación $V$. Entonces el valor obtenido es:

\begin{equation*}
    I_{MAX} = \frac{S_T}{V} = \frac{2320.148}{127} = 18.269 A
\end{equation*}
De tal manera que si todos los elementos del sistema están encendidos al mismo tiempo, esta será la corriente consumida del sistema.

\textbf{Cálculo por cada proceso del sistema}

Para esta parte, se hace el cálculo de cada uno de los procesos por separado, considerando únicamente, los subsistemas que trabajan al mismo tiempo en determinada parte del proceso.

\textit{Proceso de fermentación}

En esta parte del proceso, únicamente se tiene en cuenta el funcionamiento del circuito electrónico, para mantener el registro de los sensores, así como el funcionamiento de la interfaz humano-máquina; también se tiene activa la fuente de 24 V. Este proceso está programado para llevarse a cabo por 2 horas y media, por lo que la energía consumida en este periodo de tiempo está expresada en la tabla. \ref{tab:fermentacion_elec}.


\begin{table}
    \centering
\caption{Consumo energético en fermentación}
\label{tab:fermentacion_elec}
    \begin{tabular}{|c|c|} \hline 
         Elemento&  Potencia consumida (KW/h)\\ \hline 
         Circuito electronico&  0.050\\ \hline 
         Fuente de 24 V&  0.300\\\hline
    \end{tabular}
    
    
\end{table}

Por lo que al realizar la suma de la potencia consumida en este periodo de tiempo, da el siguiente valor:

\begin{equation*}
    P_{fer} = 0.350   KWh
\end{equation*}

\textit{Proceso de aireación}

Para este proceso, se requiere del subsistema de ventilación y el sistema de rotación del eje de aireación, también del sistema electrónico principal. Esta tarea se va a ejecutar por un ciclo de 10 minutos, por lo que el consumo de energía en esta parte del proceso está descrito en la tabla. \ref{tab:Airo_compo}.

\begin{table}
    \centering
\caption{Consumo energético en aireación}
\label{tab:Airo_compo}
    \begin{tabular}{|c|c|} \hline 
         Elemento&  Potencia consumida (KWh)\\ \hline 
         Circuito electrónico&  0.003\\ \hline 
 Motor de aireación&0.030\\ \hline 
 Ventilador&0.002\\ \hline 
         Fuente de 24 V&  0.020\\\hline
    \end{tabular}
    
    
\end{table}

Por lo que al realizar la suma de la potencia consumida en este periodo de tiempo, da el siguiente valor.

\begin{equation*}
    P_{fer} = 0.0.055   KWh
\end{equation*}

\textit{Proceso de regulación de temperatura}

Para esta parte, es necesario tener encendida la resistencia eléctrica, así como el eje de aireación y los sistemas electrónicos del sistema. Este proceso se lleva a cabo en un tiempo promedio de 20 minutos por cada hora de proceso, dependiendo de la temperatura ambiental y la pérdida de calor del contenedor. El consumo energético de este proceso está descrito en la tabla. \ref{tab:Temp_elec}.

\begin{table}
    \centering
\caption{Consumo energético cuando se eleva la temperatura}
\label{tab:Temp_elec}
    \begin{tabular}{|c|c|} \hline 
         Elemento&  Potencia consumida (KW/h)\\ \hline 
         Circuito electrónico&  0.008\\ \hline 
 Motor aireación&0.075\\ \hline 
 Resistencia eléctrica&0.417\\ \hline 
         Fuente de 24 V&  0.050\\\hline
    \end{tabular}
    
    
\end{table}

Por lo que al realizar la suma de la potencia consumida en este ciclo, se tiene el siguiente resultado.

\begin{equation*}
    P_{fer} = 0.550   KWh
\end{equation*}

\textit{Proceso de trituración}

Para este proceso, es necesario tener encendido el motor de trituración y la electrónica del sistema. Para llevar a cabo la molienda de 15 kg de materia, fue necesario mantener encendido el sistema por una hora y media, por lo que el consumo energético de este proceso está descrito en la tabla. \ref{tab:Tritu_cons}.

\begin{table}
    \centering
\caption{Consumo energético del sistema de trituración}
\label{tab:Tritu_cons}
    \begin{tabular}{|c|c|} \hline 
         Elemento&  Potencia consumida (KW/h)\\ \hline 
         Motor de trituración&  2.096\\ \hline 
         Circuito electrónico&0.040\\ \hline 
         Fuente de 24 V&  0.240\\\hline
    \end{tabular}
    
    
\end{table}

Por lo que al realizar la suma de la potencia consumida en este ciclo, se tiene el siguiente resultado.

\begin{equation*}
    P_{fer} = 2.376   KWh
\end{equation*}

\textbf{Cálculo por el total del proceso}

Para esta parte, se hace el cálculo estimado del consumo de energía eléctrica por un ciclo de 72 horas, donde los datos que se toman en cuenta, están resumidos en la tabla \ref{tab:Electrical_values_total}

\begin{table}
    \centering
\caption{Consumo total de los elementos eléctricos del sistema}
\label{tab:Electrical_values_total}
    \begin{tabular}{|m{2.5cm}|m{2.5cm}|m{2.5cm}|m{2.5cm}|} \hline 
         Elemento&  Potencia activa nominal (W)&Tiempo encendido (hr)&Energia consumid (KWh)\\ \hline 
         Motor de trituración&  1397&1.5&2.096\\ \hline 
         Motor de aireación&  180&24.5&4.410\\ \hline 
         Circuito electrónico&  20&72&1.440\\ \hline 
         Resistencia eléctrica&  1000&20&20\\ \hline 
         Ventilador&  12&6&0.072\\ \hline 
         Fuente 24 V&  120&72&8.640\\ \hline
    \end{tabular}
    
    
\end{table}

Por lo que sumando todos los valores de la ultima columna, el consumo total de energía del dispositivo es representado de la siguiente forma.

\begin{equation*}
    P_{total} = 36.658 KWh
\end{equation*}

Tomando en cuenta el costo de la luz en tarifa domestica de Enero 2024, emitido por la Comisión Federal de Electricidad, donde el kilowatt de energía cuesta \$ 1.005 MXN, el costo de la energía consumida en el proceso total es de \$ 36.840 MXN.

\clearpage